% ============================================================
% SECTION 8 — v16 DRAFT (replaces v15.2 sec:official)
% Status: DRAFT for author review.
% All numbers verified against parsed final leaderboard (n=310, July 2026)
% ============================================================

\section{Official Benchmark Validation}
\label{sec:official}

The complete competition leaderboard provides an independent validation
dataset. We use it only to evaluate whether the empirical patterns
identified in our controlled ablation study
(Sections~\ref{sec:results}--\ref{sec:multimodel}) generalize beyond our
own prompts. The analytical sequence is deliberately unidirectional: our
experiments produced the findings; the findings motivated the hypotheses
of Section~\ref{sec:analysis}; the leaderboard, released after our
submission was locked, then serves as an out-of-sample test of whether
those hypotheses are compatible with the behavior of 309 independently
developed submissions.

\subsection{Data and Scope}
\label{sec:official-data}

The final competition leaderboard~\citep{sair2026leaderboard} reports
results for 310~evaluated submissions of 1{,}007 registered participants,
across two evaluation tracks. The \emph{Evaluation (Scored)} board, which
determined official rankings, averages accuracy over three models
(GPT-OSS~120B, Gemma~4~31B~IT, Llama~3.3~70B~Instruct) and three problem
sets (\texttt{normal}, \texttt{hard}, \texttt{extra\_hard}, 200~problems
each, three runs per problem). The \emph{Order-5 (Research)}
board~\citep{sair2026leaderboardresearch} evaluates the same three models
on a separate 200-problem set restricted to order-5
magmas.\footnote{Because the two boards draw from different problem
distributions, the pair (scored accuracy, research accuracy) provides,
for every submission, exactly the kind of cross-distribution measurement
that our AN45c/AN38 case study performed for our own prompts.}

Because both boards were published after the submission deadline, no
participant---ourselves included---could optimize against them. The
leaderboard therefore functions as a natural pre-registered test: every
hypothesis stated in Sections~\ref{sec:analysis}
and~\ref{sec:multimodel} of the versions of this paper published before
competition close (v13--v15.2, timestamped on arXiv and Zenodo) predates
the evidence examined here.

Summary statistics frame the difficulty of the task. On the scored
board, mean accuracy across all 310~submissions is 55.3\% (median
55.8\%), and the maximum is 69.1\%. On the research board, the mean is
61.3\% and the maximum 82.4\%. Only 32 of 310~submissions exceed our
locally measured no-cheatsheet baseline of 59.75\% on the scored board.

\subsection{Cross-Distribution Trade-Off at Full Scale}
\label{sec:official-tradeoff}

Version~15.2 of this paper, using the voluntary Contributor Network
snapshot ($n{=}52$, 13 with dual-split scores), observed that only one
submission exceeded 65\% accuracy on both hard2 and hard3, and flagged
the full leaderboard as the natural falsification test: the
cross-distribution trade-off surface would either survive at scale or
the thesis would fail. The complete leaderboard resolves this test.

\begin{table}[t]
\centering
\caption{Cross-board performance thresholds, final leaderboard
($n{=}310$). Each submission is evaluated on both the official scored
distribution and the order-5 research distribution.}
\label{tab:dualthreshold}
\begin{tabular}{lc}
\toprule
Criterion & Submissions \\
\midrule
$>$60\% on both boards & 27 / 310 \\
$>$65\% on both boards & 5 / 310 \\
$>$70\% on both boards & \textbf{0 / 310} \\
\bottomrule
\end{tabular}
\end{table}

No submission among all 310 exceeds 70\% accuracy on both distributions
(Table~\ref{tab:dualthreshold}). The trade-off is most visible at the
top of the research board: the five highest-scoring research submissions
lose between 19 and 32~percentage points when evaluated on the official
scored distribution (Table~\ref{tab:researchcollapse}), a pattern
structurally identical to the AN45c inversion documented in
Section~\ref{sec:official-inversion} (79.25\% local $\rightarrow$ 55.5\%
official hard3, a 23.75pp drop).

\begin{table}[t]
\centering
\caption{The five highest-ranked submissions on the research board and
their official scored results. Team names are public leaderboard
entries~\citep{sair2026leaderboard, sair2026leaderboardresearch}.}
\label{tab:researchcollapse}
\begin{tabular}{lccc}
\toprule
Team & Research & Scored & Scored rank \\
\midrule
Weihang Ding & 82.4\% & 56.3\% & 137 \\
Lothar Schiemanowski & 81.9\% & 50.3\% & 288 \\
Math Distillation & 80.6\% & 61.6\% & 20 \\
Ken & 80.4\% & 56.2\% & 141 \\
zkdmt & 80.1\% & 52.2\% & 264 \\
\bottomrule
\end{tabular}
\end{table}

The exception identified in v15.2 also resolves. Arjun Garg's
submission, which exceeded 65\% on both hard2 and hard3 in the
Contributor Network snapshot and which we flagged as a potential
falsifier of the trade-off surface, ranks 8th on the research board
(77.9\%) but 166th on the scored board (55.5\%)---below our own
submission's scored accuracy. What appeared at $n{=}13$ to be an escape
from the trade-off surface is, at $n{=}310$, another instance of it.

Conversely, the five submissions that do clear 65\% on both boards
(Table~\ref{tab:robustfive}) are characterized not by peak accuracy on
either distribution but by small cross-board gaps. The competition
winner exhibits the smallest gap of the top group
($+1.0$pp), despite ranking only 15th--20th on the research board. At
full scale, cross-distribution robustness---not peak accuracy on any
single distribution---predicts official ranking. This generalizes,
across 309 independently developed prompts, the lesson of our own
AN38-over-AN45c submission decision (Section~\ref{sec:results}).

\begin{table}[t]
\centering
\caption{All submissions exceeding 65\% on both boards ($n{=}5$ of
310), ordered by scored accuracy.}
\label{tab:robustfive}
\begin{tabular}{lccc}
\toprule
Team & Scored & Research & Gap \\
\midrule
Dufius & 69.1\% & 70.1\% & $+1.0$ \\
vt & 67.9\% & 66.2\% & $-1.7$ \\
Reza Jamei & 66.3\% & 74.1\% & $+7.8$ \\
AJ & 66.2\% & 70.2\% & $+4.0$ \\
yan-biao & 65.3\% & 74.3\% & $+9.0$ \\
\bottomrule
\end{tabular}
\end{table}

Two conclusions follow. First, the cross-distribution fragility
documented for our own prompts in
Section~\ref{sec:official-inversion} is not an idiosyncrasy of our
design family: it is the dominant failure mode of the entire field of
310~submissions. Second, the empirical saturation region identified in
Section~\ref{sec:ceiling} for our prompt family (approximately
60--79\% on local hard3) is consistent with the official ceiling
observed across all participants: no submission exceeds 69.1\% on the
official scored evaluation, and no submission exceeds 70\% on both
distributions simultaneously.

\subsection{Compatibility with the Structural Classification Hypothesis}
\label{sec:official-router}

Section~\ref{sec:analysis} advanced the hypothesis---prior to the
release of any official results---that gpt-oss-120b behaves on this
task as a structural-pattern classifier rather than a prover, and that
prompts succeed to the extent that they support structural
classification rather than open-ended mathematical reasoning. The
leaderboard offers an observational test: if the hypothesis is
directionally correct, strategies that explicitly constrain the model's
decision process should be over-represented at the top of the official
ranking.

The evidence is compatible with the hypothesis. The winning submission
is, by its own publicly visible description, an explicitly rule-ordered
deterministic classifier: a fixed sequence of syntactic rules applied
in priority order with a first-match stopping condition, rather than a
reasoning guide.\footnote{We cite only the publicly visible
categorical structure of leaderboard submissions and their published
scores. A content-level analysis of other participants' prompts is
outside the scope of this paper.} We treat this as external supporting
evidence, not as proof: a single winning instance cannot establish
that decision-space reduction is necessary for escaping the saturation
region, and the leaderboard does not permit controlled comparison of
strategy categories. What the observation does establish is that the
strategy our hypothesis predicted would outperform reasoning
guidance---formulated in Section~\ref{sec:analysis} before official
results existed---is the strategy that won.

Section~\ref{sec:instability} develops the observable quantity that,
in our view, underlies this pattern: the degree to which a prompt
delegates the classification decision to the model, measurable through
run-to-run verdict instability.
